\section{Security Overview}

Security is treated as an integral part of the SypIT Homelab architecture
rather than as a separate component added after infrastructure deployment.

The environment is designed according to the principle of defence in depth.
No single security mechanism is considered sufficient to protect the entire
infrastructure. Instead, multiple independent security controls will be
implemented across network, system, application, identity, and data layers.

The security architecture will evolve together with the infrastructure.
Controls described in this chapter therefore represent both currently
implemented mechanisms and security requirements for future development
phases.


\section{Security Principles}

The project follows several fundamental security principles.

\subsection{Least Privilege}

Users, applications, and services should receive only the permissions
required to perform their intended functions.

Administrative privileges should not be used for regular operations when
lower privileges are sufficient.

\subsection{Deny by Default}

Network communication and access to resources should be denied unless
explicitly permitted.

This principle will be applied particularly to firewall rules,
network segmentation, database access, and service authorization.

\subsection{Defence in Depth}

Security controls should exist at multiple independent layers.

A failure or misconfiguration of one control should therefore not
automatically expose the entire infrastructure.

The planned security layers include:

\begin{itemize}
    \item perimeter firewalling,
    \item network segmentation,
    \item host-based firewalls,
    \item operating system hardening,
    \item identity and access management,
    \item application-level authorization,
    \item encrypted communication,
    \item monitoring and logging,
    \item backup and recovery mechanisms.
\end{itemize}

\subsection{Separation of Duties}

Infrastructure components with different security requirements should
be separated whenever practical.

For example, publicly accessible services should not share unrestricted
access with management systems or database infrastructure.


\section{Network Security}

Network security is one of the primary security layers of the project.

The initial laboratory network is intentionally simple. As the environment
grows, the network architecture will be divided into multiple logical
security zones.

Planned zones include:

\begin{itemize}
    \item Management,
    \item Servers,
    \item Databases,
    \item IoT and Robotics,
    \item Game Servers,
    \item DMZ,
    \item Client Devices.
\end{itemize}

Communication between these zones will be controlled by routing and
firewall policies.

The target communication model follows the principle:

\begin{center}
\texttt{deny by default, allow by exception}
\end{center}

For example, a web application located in the DMZ may be permitted to
communicate with a specific backend service or database port, while direct
access from the DMZ to the management network should remain prohibited.


\section{DMZ and Public Services}

Services accessible from external networks represent a higher security risk
than internal services.

Public-facing workloads may include:

\begin{itemize}
    \item websites,
    \item web applications,
    \item APIs,
    \item file transfer services,
    \item game servers.
\end{itemize}

These services should eventually operate inside dedicated network segments
or a DMZ.

Compromise of a public service should not provide unrestricted access to
internal infrastructure.
\newpage
The intended model is:

\begin{verbatim}
                    Internet
                       |
                 Router/Firewall
                       |
                      DMZ
                 /      |      \
               WWW    Games    File
                |
          controlled access
                |
          Internal Services
\end{verbatim}


\section{Administrative Access}

Administrative access should be separated from normal application traffic.

Remote administration of Linux systems will primarily use SSH.

Password-based remote authentication should be minimized where practical
in favour of public-key authentication.

Administrative access from external networks should eventually be provided
through a dedicated VPN rather than by directly exposing administrative
services to the Internet.

The target model is:

\begin{verbatim}
Administrator
      |
     VPN
      |
Management Network
      |
     SSH
      |
Managed Systems
\end{verbatim}

Direct exposure of SSH, database administration interfaces, hypervisor
management interfaces, or monitoring administration panels to the public
Internet should be avoided.


\section{Identity and Access Management}

Access to infrastructure resources should be associated with identifiable
users or services.

Shared administrative accounts should be avoided where practical.

Administrative activities should use individual accounts and privilege
escalation mechanisms such as \texttt{sudo}.

Service accounts should be separated from interactive user accounts and
should receive only the permissions necessary for their workloads.

Future development may introduce centralized identity management for
selected infrastructure components.


\section{Secrets Management}

Sensitive information must not be stored directly in source code or
committed to the Git repository.

Examples of sensitive information include:

\begin{itemize}
    \item passwords,
    \item API tokens,
    \item database credentials,
    \item private SSH keys,
    \item private TLS keys,
    \item application secrets,
    \item authentication tokens.
\end{itemize}

The repository uses \texttt{.gitignore} rules to reduce the risk of
accidentally committing common files containing secrets.

However, \texttt{.gitignore} is not considered a security mechanism by
itself.

Future phases of the project will introduce centralized secrets management
and controlled secret distribution to applications and infrastructure
components.


\section{Encryption and PKI}

Sensitive communication should use encrypted protocols whenever possible.

Examples include:

\begin{itemize}
    \item HTTPS instead of HTTP for sensitive or externally accessible traffic,
    \item SSH instead of insecure remote shell protocols,
    \item SFTP or FTPS where secure file transfer is required,
    \item encrypted database connections,
    \item VPN tunnels for remote administrative access.
\end{itemize}

The project is expected to introduce an internal Public Key Infrastructure
(PKI) for laboratory services.

This will allow practical implementation of:

\begin{itemize}
    \item Certificate Authorities,
    \item server certificates,
    \item client certificates,
    \item TLS,
    \item mutual TLS,
    \item certificate renewal,
    \item certificate revocation and rotation.
\end{itemize}


\section{Operating System Security}

Linux systems should be maintained according to a consistent baseline.

The security baseline will include:

\begin{itemize}
    \item regular security updates,
    \item minimal installation of unnecessary services,
    \item controlled administrative privileges,
    \item host-based firewall configuration,
    \item secure SSH configuration,
    \item appropriate file permissions,
    \item service isolation,
    \item system logging,
    \item periodic configuration review.
\end{itemize}

Server workloads should eventually use dedicated server-oriented operating
systems rather than relying exclusively on the initial Ubuntu Desktop
development machine.


\section{Monitoring and Security Visibility}

Security controls must provide sufficient visibility to detect abnormal
behaviour.

The project will therefore include centralized collection and analysis of
system, application, and network information.

Planned capabilities include:

\begin{itemize}
    \item centralized logging,
    \item infrastructure metrics,
    \item application monitoring,
    \item network traffic analysis,
    \item security alerts,
    \item IDS/IPS experimentation,
    \item audit information.
\end{itemize}

Monitoring should support both operational troubleshooting and security
incident investigation.

\newpage
\section{Backup and Recovery}

Backups are considered part of the security architecture.

Critical configuration, databases, application data, and infrastructure
state should be backed up according to defined policies.

A backup is not considered reliable until its restoration procedure has
been successfully tested.

Future stages will therefore define:

\begin{itemize}
    \item backup scope,
    \item backup frequency,
    \item retention policies,
    \item storage locations,
    \item recovery procedures,
    \item Recovery Point Objectives (RPO),
    \item Recovery Time Objectives (RTO).
\end{itemize}


\section{Security Testing}

The laboratory environment provides an isolated platform for controlled
security testing.

Planned experiments may include:

\begin{itemize}
    \item firewall rule validation,
    \item network segmentation tests,
    \item vulnerability scanning,
    \item service exposure verification,
    \item authentication testing,
    \item certificate expiration scenarios,
    \item network failure simulations,
    \item intrusion detection experiments.
\end{itemize}

Testing should be performed only against systems belonging to the laboratory
environment or systems for which explicit authorization has been granted.


\section{Incident Management}

Operational and security incidents should be documented in a structured
manner.

For significant incidents, the following information should be recorded:

\begin{itemize}
    \item incident description,
    \item detection method,
    \item affected components,
    \item technical impact,
    \item root cause,
    \item remediation,
    \item preventive actions,
    \item lessons learned.
\end{itemize}

This process will also be used during future chaos engineering exercises,
where selected infrastructure failures will be intentionally introduced
to evaluate monitoring, diagnosis, and recovery procedures.


\section{Security Development Strategy}

Security mechanisms will be introduced incrementally as the infrastructure
grows.

The planned progression is:

\begin{enumerate}
    \item secure the initial Linux and virtualization environment,
    \item establish controlled administrative access,
    \item introduce network segmentation,
    \item deploy firewall policies,
    \item isolate public-facing services,
    \item introduce VPN-based remote administration,
    \item implement centralized secrets management,
    \item deploy internal PKI and encrypted service communication,
    \item centralize logs and security monitoring,
    \item introduce IDS/IPS capabilities,
    \item implement and test backup and disaster recovery,
    \item perform controlled failure and security testing.
\end{enumerate}

The security architecture will be reviewed whenever new services,
network segments, or externally accessible components are introduced.